\documentclass[twocolumn,3p,times,procedia]{elsarticle}

\usepackage{ecrc}
\usepackage{graphicx}
\usepackage{balance}
\usepackage{array,threeparttable,tabularx,booktabs,multirow,makecell}
\usepackage{url}
\usepackage{amssymb,amsmath,amsthm,bm}
\usepackage[figuresright]{rotating}
\usepackage{caption}
\usepackage{titlesec}
\usepackage{fancyhdr}
\usepackage{geometry}
\usepackage{xcolor}
\usepackage{float}
\usepackage{placeins}
\usepackage{fontspec}
\usepackage{xeCJK}
\usepackage[hidelinks]{hyperref}
\usepackage{tikz}
\usetikzlibrary{arrows.meta,positioning,fit,calc,shapes.geometric}

\IfFontExistsTF{SimSun}{\setCJKmainfont{SimSun}}{\setCJKmainfont{FandolSong-Regular}}
\IfFontExistsTF{SimHei}{\setCJKsansfont{SimHei}}{\setCJKsansfont{FandolHei-Regular}}
\IfFontExistsTF{Times New Roman}{\setmainfont{Times New Roman}}{}

\volume{00}
\firstpage{1}
\runauth{作者信息待补}
\jid{procs}
\CopyrightLine{2026}{Published by Elsevier Ltd.}

\captionsetup[figure]{labelfont={footnotesize,bf},name={图},labelsep=period,font={footnotesize}}
\captionsetup[table]{labelfont={footnotesize,bf},name={表},labelsep=newline,singlelinecheck=false,font={footnotesize}}
\makeatletter
\setlength{\@fptop}{0pt}
\setlength{\@fpsep}{10pt}
\setlength{\@fpbot}{0pt plus 1fil}
\setlength{\@dblfptop}{0pt}
\setlength{\@dblfpsep}{10pt}
\setlength{\@dblfpbot}{0pt plus 1fil}
\makeatother
\titleformat*{\section}{\small \bfseries}
\titleformat*{\subsection}{\small \em}
\titleformat*{\subsubsection}{\small \em}
\geometry{left=1.25cm,right=1.15cm,top=1.9cm,bottom=1.9cm,foot=1.05cm}
\setlength\columnsep{0.6cm}
\pagestyle{fancy}
\fancyhf{}
\fancyheadoffset[RO,EL]{0pt}
\fancyhead[RO,LE]{\footnotesize \thepage}
\fancyhead[ER]{\em \footnotesize 作者信息待补}
\fancyhead[LO]{\em \footnotesize Network-Resilient Cooperative Transport Control}

\newcommand{\todozh}[1]{\textcolor{red}{[待补：#1]}}
\newcommand{\AKE}{AKE-M}
\newcommand{\Method}{NR-KDCC}
\newcommand{\MethodCN}{网络韧性 Koopman-时延一致性协同运输控制}
\newcommand{\TF}{\Method{}}
\newcommand{\R}{\mathbb{R}}
\newcommand{\norm}[1]{\left\lVert #1 \right\rVert}
\renewcommand{\refname}{参考文献}
\makeatletter
\renewenvironment{abstract}{\global\setbox\absbox=\vbox\bgroup
  \hsize=\textwidth\def\baselinestretch{1}%
  \noindent\unskip\textbf{摘要}
 \par\medskip\noindent\unskip\ignorespaces}
 {\egroup}
\def\keyword{%
  \def\sep{\unskip, }%
  \def\MSC{\@ifnextchar[{\@MSC}{\@MSC[2000]}}%
  \def\@MSC[##1]{\par\leavevmode\hbox {\it ##1~MSC:\space}}%
  \def\PACS{\par\leavevmode\hbox {\it PACS:\space}}%
  \def\JEL{\par\leavevmode\hbox {\it JEL:\space}}%
  \global\setbox\keybox=\vbox\bgroup\hsize=\textwidth
  \normalsize\normalfont\def\baselinestretch{1}
  \parskip\z@
  \noindent\textit{关键词：}
  \raggedright
  \ignorespaces}
\makeatother
\newtheorem{assumption}{假设}
\newtheorem{theorem}{定理}
\newtheorem{remark}{注}
\definecolor{nrInk}{HTML}{243B53}
\definecolor{nrBlue}{HTML}{2F5D8C}
\definecolor{nrTeal}{HTML}{008C8C}
\definecolor{nrOrange}{HTML}{D97732}
\definecolor{nrRose}{HTML}{B6465F}
\definecolor{nrSlate}{HTML}{5E7184}
\definecolor{nrPaleBlue}{HTML}{EAF2FA}
\definecolor{nrPaleGreen}{HTML}{E8F6F3}
\definecolor{nrPaleOrange}{HTML}{FFF1E3}
\tikzset{
  nrbox/.style={rounded corners=2pt, draw=nrBlue, very thick, fill=nrPaleBlue, align=center, inner sep=4pt, font=\footnotesize},
  nrmodule/.style={rounded corners=2pt, draw=nrTeal, thick, fill=nrPaleGreen, align=center, inner sep=4pt, font=\footnotesize},
  nrsafe/.style={rounded corners=2pt, draw=nrRose, thick, fill=nrPaleOrange, align=center, inner sep=4pt, font=\footnotesize},
  nrarr/.style={-{Latex[length=2.2mm]}, thick, draw=nrSlate},
  nrdash/.style={-{Latex[length=2.2mm]}, thick, dashed, draw=nrRose}
}

\begin{document}\small
\begin{frontmatter}

\dochead{}
\title{
\begin{flushleft}
{\LARGE 基于双线性 Koopman 学习与时延补偿一致性 MPC 的网络韧性协同运输控制}
\end{flushleft}
}

\author[]{\leftline{作者信息待补$^{a,*}$}}
\address{\leftline{$^a$单位信息待补，城市，邮编，国家}}

\begin{abstract}
多车协同运输在仓储物流、园区配送和柔性制造中具有明确应用价值，但其闭环控制同时受到车辆非线性动力学、共享载荷几何约束、通信时延/丢包和执行器退化的耦合影响。已有自适应 Koopman 控制方法可为复杂系统提供数据驱动预测模型，但通常不显式处理多车刚性载荷连接安全、网络退化一致性约束和故障降级保护。本文在上传基线工作的思想基础上，构建面向四车刚性载荷运输的 \Method{} 控制框架：首先在大地坐标系与路径坐标系下建立车辆--载荷耦合模型，并离线生成覆盖 nominal、通信退化、单车故障与混合扰动的数据；随后采用稳定投影双线性 Koopman 提升模型描述控制输入与非线性动力学的耦合；最后在 MPC 中加入通信质量感知一致性、时延状态外推、约束收缩、执行器效率估计、FTC 重分配和证书触发 fallback。当前多随机种子实验表明，在 mixed fault/noise 场景下，相对任务对齐的 \AKE{} 机制基线，\TF{} 将横向 RMSE 从 0.0453 降至 0.0408，将连接最大利用率从 0.2999 降至 0.0513；在 high-communication-noise 场景下，横向 RMSE 从 0.0519 降至 0.0424，连接最大利用率从 0.3757 降至 0.0452。与此同时，力峰值出现上升，因此本文将载荷受力作为必须审计的代价项而非单调改善指标。本文给出输入到状态实践有界意义下的 Lyapunov 证明，并把 Koopman 模块限定为控制窗口内可审计预测结构，而不宣称远期开环 rollout 全面优于线性模型。
\end{abstract}

\begin{keyword}
协同运输 \sep Koopman 算子 \sep 双线性系统 \sep 模型预测控制 \sep 网络时延 \sep 丢包 \sep 容错控制 \sep Lyapunov 稳定性
\end{keyword}

\end{frontmatter}

\section{引言}

移动机器人与智能车辆的协同运输任务要求多个执行体在保持载荷几何关系的同时完成轨迹跟踪。与单车路径跟踪相比，该问题具有三类额外困难：第一，共享载荷将各车辆的横向、纵向和姿态误差耦合在一起；第二，车辆间通信质量会直接影响一致性目标和载荷连接误差；第三，执行器效率下降或控制饱和会通过连接点传递为载荷受力波动。若只报告单车轨迹误差或 formation error，难以说明运输过程是否安全。

Koopman 学习为非线性车辆控制提供了可优化预测模型。近年 Koopman-MPC 已用于自动驾驶车辆横向控制、随机 MPC 和深度 Koopman 表示学习，例如 K-SMPC 将 Koopman 预测误差纳入 stochastic MPC，深度 Koopman-ESO 方法将观测器用于补偿车辆模型残差，核 Koopman 误差界研究也强调有限数据预测器需要误差边界和保守保护\cite{kim2025ksmpc,chen2024esoDeepKoopman,philipp2023koopmanErrorBounds,cibulka2021koopmanVehicleMpc}。这些研究说明 Koopman 模型适合为 MPC 提供局部线性或近线性的预测结构，但主要面向单车轨迹/车道保持，不直接处理四车刚性载荷运输中的连接误差、载荷受力和通信退化耦合。

多机器人协同运输研究已经发展出 formation control、分布式优化、非完整约束处理和学习型协作策略。近期差速移动机器人协同搬运、嵌入式 formation 搬运、多目标强化学习运输等工作分别从控制结构、平台验证和任务决策层面推进了该方向\cite{cooperativeTransport2024ras,embeddedFormationTransport2024cep,multiObjectiveTransport2024mechatronics,intelligentCollaborativeTransport2024}。然而，很多协同运输工作将通信视为理想条件或实现细节，缺少对 packet loss、delay 和 link quality 如何影响载荷连接安全与受力审计的闭环分析。

网络化 MPC 和 CAV/DMPC 文献近三年开始系统研究通信退化。已有工作考虑 Markov packet loss、DoS attack、communication loss、随机时延和网络延迟缓解策略\cite{bian2025markovPacketLossDMPC,dosDMPC2024isa,robustCommunicationLoss2025trb,bao2024robustUkfMpc,networkLatencyCav2024sensors,dai2024networkedPredictiveControl}。这些方法多服务于 longitudinal platoon、spacing regulation 或 string stability。协同运输的关键安全变量并不是车距串稳定性，而是载荷中心误差、角点连接误差、车辆输入约束和载荷受力。因此，通信退化需要进入协同运输 MPC 的一致性权重、邻居预测、约束收缩和 fallback 触发逻辑。

容错与鲁棒 MPC 方面，tube/convex robust DMPC、learning-supported MPC 和网络化预测控制为处理模型不确定性与外部扰动提供了基础\cite{convexRobustDMPC2025ejc,gasparino2023nnMpcUncertainty,dai2024networkedPredictiveControl}。但在本文当前证据链中，执行器 FDI/FTC 更适合作为闭环安全增强模块，而不是独立宣称超越现有 FTC 文献的主理论贡献。原因是现阶段仍需补 1--2 篇近三年车辆/多机器人 actuator fault-tolerant MPC 或 FDI/control allocation 的直接文献，并补充故障效率估计曲线与重分配过程图。

基于上述缺口，本文提出 \MethodCN{}（\Method{}）框架。该命名对应本文的四个核心对象：网络韧性、Koopman 双线性预测、时延补偿一致性 MPC 和协同运输闭环。与上传基线 ``Adaptive Koopman Embedding for Robust Control of Complex'' 的核心思想相比，本文不只使用自适应 Koopman 预测器，还将学习模型、通信质量、载荷连接安全和故障降级共同放入可审计闭环。需要强调的是，当前实验中的主基线是任务对齐的 \AKE{} 机制基线，而不是严格复现上传原文全部设置的 AKE-A；因此本文只在当前四车载荷运输场景中报告相对 \AKE{} 的改善，不写成对原文 AKE-A 的全面超越。

本文贡献如下。

\begin{enumerate}
\item 在 Koopman 车辆 MPC 研究基础上，提出稳定投影双线性 Koopman 预测器，并将其嵌入四车刚性载荷协同运输 MPC。其机制是保留 lifted dynamics 中控制输入对状态演化的乘性影响，同时通过谱投影和证书项限制长时域预测漂移。
\item 在 CAV/DMPC 通信退化研究基础上，提出通信质量感知的时延补偿一致性 MPC。其机制是将丢包、时延和链路质量映射为邻居状态外推、一致性权重、输入约束收缩和 degraded fallback，从而直接保护载荷连接误差。
\item 在协同运输安全评价基础上，建立载荷中心、角点连接误差、车辆输入和载荷受力的联合审计链。其机制是将运输性能和物理安全拆成不同指标报告，避免用单一轨迹误差掩盖连接风险或受力代价。
\item 集成执行器效率估计、FTC 控制重分配和 Lyapunov/ISS 证书保护。其机制是在故障残差或通信质量触发阈值时，从性能优先切换为连接/受力安全优先的保守模式；该贡献当前作为安全增强模块表述。
\end{enumerate}

\section{系统建模与问题定义}

\subsection{大地坐标系车辆动力学}

考虑 $N=4$ 辆车辆协同搬运刚性载荷。第 $i$ 辆车在大地坐标系下的位置和航向为
$p_i=[X_i,Y_i]^\top$、$\psi_i$，车体系纵向/横向速度和横摆角速度为 $u_i$、$v_i$、$r_i$。控制输入为纵向加速度 $a_{x,i}$ 和前轮转角 $\delta_i$。车体系到大地坐标系的运动学为
\begin{equation}
\begin{aligned}
\dot X_i &= u_i\cos\psi_i-v_i\sin\psi_i,\\
\dot Y_i &= u_i\sin\psi_i+v_i\cos\psi_i,\\
\dot \psi_i &= r_i .
\end{aligned}
\label{eq:world_kinematics}
\end{equation}
采用线性轮胎侧偏模型时，车辆侧向和横摆动力学可写为
\begin{equation}
\begin{aligned}
\dot v_i=&-\frac{2C_f+2C_r}{m u_i}v_i+
\left(-u_i-\frac{2C_f l_f-2C_r l_r}{m u_i}\right)r_i
+\frac{2C_f}{m}\delta_i+d_{v,i},\\
\dot r_i=&-\frac{2C_f l_f-2C_r l_r}{I_z u_i}v_i
-\frac{2C_f l_f^2+2C_r l_r^2}{I_z u_i}r_i
+\frac{2C_f l_f}{I_z}\delta_i+d_{r,i},\\
\dot u_i=&a_{x,i}+r_i v_i+d_{u,i}.
\end{aligned}
\label{eq:bicycle_dynamics}
\end{equation}
其中 $m,I_z,l_f,l_r,C_f,C_r$ 分别为车辆质量、横摆转动惯量、前/后轴距和前/后轮侧偏刚度；$d_{u,i},d_{v,i},d_{r,i}$ 表示模型误差、载荷耦合扰动或外部扰动。为避免低速奇异，仿真中使用 $u_i\leftarrow\max(u_i,u_{\min})$，当前代码取 $u_{\min}=0.5$。

\subsection{路径坐标误差}

设参考路径弧长为 $s$，曲率为 $\kappa(s)$，路径切向航向为 $\psi_r(s)$。车辆投影到路径坐标系后，横向误差为 $e_{y,i}$，航向误差为 $e_{\psi,i}=\psi_i-\psi_r(s_i)$。标准 Frenet 误差动力学为
\begin{equation}
\dot s_i=\frac{u_i\cos e_{\psi,i}-v_i\sin e_{\psi,i}}
{1-\kappa(s_i)e_{y,i}},
\label{eq:s_dot}
\end{equation}
\begin{equation}
\dot e_{y,i}=u_i\sin e_{\psi,i}+v_i\cos e_{\psi,i},\qquad
\dot e_{\psi,i}=r_i-\kappa(s_i)\dot s_i .
\label{eq:frenet_error}
\end{equation}
式 \eqref{eq:world_kinematics} 用于大地坐标系仿真和载荷几何约束，式 \eqref{eq:s_dot}--\eqref{eq:frenet_error} 用于构造轨迹跟踪误差和 MPC 代价。

\subsection{刚性载荷与连接误差}

载荷中心位姿为 $q_L=[X_L,Y_L,\psi_L]^\top$。第 $i$ 个连接点在载荷坐标系中的固定向量为 $b_i$，对应大地坐标位置为
\begin{equation}
p_{L,i}=p_L+R(\psi_L)b_i,\qquad
R(\psi)=\begin{bmatrix}\cos\psi&-\sin\psi\\ \sin\psi&\cos\psi\end{bmatrix}.
\label{eq:payload_anchor}
\end{equation}
车辆与载荷连接误差定义为
\begin{equation}
e_{c,i}=p_i-p_{L,i},\qquad
e_c=\mathrm{col}(e_{c,1},\ldots,e_{c,N}).
\label{eq:connection_error}
\end{equation}
若采用等效弹簧阻尼连接，则连接力可表示为
\begin{equation}
F_{c,i}=K_c e_{c,i}+D_c\dot e_{c,i},
\label{eq:connection_force}
\end{equation}
并通过 $\sum_i F_{c,i}$ 与 $\sum_i (R(\psi_L)b_i)\times F_{c,i}$ 作用到载荷平动和转动。本文在控制层不声称连接力峰值单调降低，而是将 $\max_i\norm{F_{c,i}}$、RMS 和方向分量作为安全代价监测指标。

\subsection{通信图与退化模型}

车辆通信拓扑记为 $\mathcal G_k=(\mathcal V,\mathcal E_k)$，边 $(j,i)$ 的链路质量为 $q_{ij,k}\in[0,1]$，测得时延为 $\tau_{ij,k}$，丢包指示为 $\ell_{ij,k}\in\{0,1\}$。当邻居状态可用时，控制器接收 $\hat x_{j,k-\tau}$；当时延或丢包发生时，使用 Koopman/运动学预测器进行外推：
\begin{equation}
\hat x_{j,k|k}^{\mathrm{net}}=
\begin{cases}
\Pi_x \hat z_{j,k-\tau+\tau|k-\tau}, & \ell_{ij,k}=0,\\
\hat x_{j,k-1|k-1}^{\mathrm{net}}, & \ell_{ij,k}=1 .
\end{cases}
\label{eq:delay_prediction}
\end{equation}
这里 $\Pi_x$ 为 lifted state 到物理状态的投影。通信质量进一步决定一致性权重和约束收缩比例，使网络退化不再只是外部扰动，而是 MPC 中的调度变量。

\section{\Method{} 方法}

\subsection{方法总览与命名}

本文方法命名为 \MethodCN{}（\Method{}），而不再使用本地实现版本号命名。其逻辑链条为：首先在大地坐标系和路径坐标系下建立车辆--载荷--通信联合模型；然后用该模型离线生成覆盖曲率、通信退化和执行器效率下降的数据；接着训练带稳定投影的双线性 Koopman 预测器；最后将预测器嵌入通信质量感知的一致性 MPC，并由 FDI/FTC、约束收缩和 Lyapunov/ISS 证书提供安全保护。这样组织后，方法模块、贡献陈述和后续实验指标一一对应：预测模型对应动力学学习和误差时间序列，通信补偿对应 high-communication-noise 对比与消融，连接/受力安全对应 E7 审计，FDI/FTC 与证书对应故障场景和稳定性统计。

\begin{figure*}[t]
\centering
\resizebox{0.98\textwidth}{!}{%
\begin{tikzpicture}[node distance=6mm and 6mm]
\tikzset{nrdata/.style={rounded corners=2pt, draw=nrOrange, thick, fill=nrPaleOrange, align=center, inner sep=4pt, font=\footnotesize}}
\node[nrbox, text width=2.75cm] (scenario) {任务与先验\\参考路径 $r(s),\kappa(s)$\\载荷几何 $b_i$\\车辆参数 $m,I_z,C_f,C_r$};
\node[nrbox, text width=3.0cm, right=of scenario] (model) {车辆--载荷建模\\大地坐标 $(X_i,Y_i,\psi_i)$\\路径误差 $(s_i,e_{y,i},e_{\psi,i})$\\连接 $e_{c,i},F_{c,i}$};
\node[nrdata, text width=2.9cm, right=of model] (offline) {离线数据生成\\曲率/输入激励\\时延、丢包、链路质量\\执行器效率 $\eta_i$\\记录 $\mathcal D$};
\node[nrmodule, text width=3.0cm, right=of offline] (koop) {双线性 Koopman 学习\\$z=[x;\phi_\theta(x)]$\\$z^+=Az+Bu+\sum u_\ell N_\ell z$\\岭回归 $\lambda$\\谱投影 $\Pi_\rho$};

\node[nrbox, text width=2.75cm, below=of scenario] (measure) {在线测量/估计\\车辆 $x_{i,k}$\\载荷 $q_{L,k}$\\上一控制 $u_{k-1}$\\残差 $r^{\mathrm{FDI}}_k$};
\node[nrmodule, text width=3.0cm, below=of model] (comm) {通信监测\\拓扑 $\mathcal G_k$\\链路质量 $q_{ij,k}$\\时延 $\tau_{ij,k}$\\丢包 $\ell_{ij,k}$\\EWMA 平滑};
\node[nrmodule, text width=2.9cm, below=of offline] (delay) {时延补偿邻居预测\\可用：Koopman 外推\\丢包：保持/短时预测\\得到 $\hat x^{\mathrm{net}}_{j,k|k}$};
\node[nrmodule, text width=3.0cm, below=of koop] (mpc) {一致性 MPC\\跟踪 $e_L,e_y,e_\psi$\\一致性 $\omega_{ij,k}\|x_i-\hat x_j\|$\\连接 $e_c$\\输入 $u,\Delta u$};

\node[nrsafe, text width=2.75cm, below=of measure] (constraints) {约束与收缩\\$u,\delta,a_x,\Delta u$\\$\|e_{c,i}\|\le e_{c,\max}(q)$\\$\|F_{c,i}\|\le F_{c,\max}$};
\node[nrsafe, text width=3.0cm, below=of comm] (ftc) {FDI/FTC 与角色调度\\估计 $\hat\eta_i$\\控制重分配\\phase-role 权重\\故障/通信退化切换};
\node[nrsafe, text width=2.9cm, below=of delay] (cert) {Lyapunov/ISS 证书\\$V_{k+1}-V_k$ 裕度\\ok ratio\\degraded fallback\\安全可行优先};
\node[nrbox, text width=3.0cm, below=of mpc] (audit) {输出与实验审计\\团队中心轨迹/误差\\各车误差\\连接利用率\\载荷受力 $F_x,F_y,M_z$\\证书统计};

\draw[nrarr] (scenario) -- (model);
\draw[nrarr] (model) -- (offline);
\draw[nrarr] (offline) -- (koop);
\draw[nrarr] (koop) -- (mpc);
\draw[nrarr] (measure) -- (comm);
\draw[nrarr] (comm) -- (delay);
\draw[nrarr] (delay) -- (mpc);
\draw[nrarr] (measure) -- (constraints);
\draw[nrarr] (constraints) -- (ftc);
\draw[nrarr] (ftc) -- (cert);
\draw[nrarr] (cert) -- (audit);
\draw[nrarr] (mpc) -- (audit);
\draw[nrarr] (constraints.east) -| (mpc.south);
\draw[nrarr] (ftc.north east) -- (mpc.south west);
\draw[nrdash] (audit.east) -- ++(0.45,0) |- node[pos=0.25, right, font=\scriptsize, color=nrRose]{闭环反馈/日志} (measure.east);
\draw[nrdash] (cert.north) -- node[right, font=\scriptsize, color=nrRose]{不安全时收缩/降级} (mpc.south);

\node[draw=nrBlue, rounded corners=2pt, fit=(scenario)(model)(offline)(koop), inner sep=5pt, label={[nrBlue]above:离线建模与学习层}] {};
\node[draw=nrTeal, rounded corners=2pt, fit=(measure)(comm)(delay)(mpc), inner sep=5pt, label={[nrTeal]left:在线预测与优化层}] {};
\node[draw=nrRose, rounded corners=2pt, fit=(constraints)(ftc)(cert)(audit), inner sep=5pt, label={[nrRose]below:安全保护与证据审计层}] {};
\end{tikzpicture}}
\caption{\Method{} 详细总体框架图。含义：该图把本文方法拆成三层。第一层从车辆参数、载荷几何和参考路径出发，建立大地坐标/路径坐标联合模型，并生成覆盖通信退化与执行器效率下降的离线数据；第二层用双线性 Koopman 模型提供在线预测，并把链路质量、时延和丢包转化为邻居状态外推和一致性 MPC 权重；第三层通过约束收缩、FDI/FTC、角色调度和 Lyapunov/ISS 证书形成安全保护，并输出团队中心误差、各车误差、连接利用率和载荷受力等实验审计指标。}
\label{fig:nrkdcc_framework}
\end{figure*}

\begin{figure*}[t]
\centering
\resizebox{0.98\textwidth}{!}{%
\begin{tikzpicture}[node distance=5.5mm and 6mm]
\tikzset{nrdecision/.style={diamond, aspect=2.2, draw=nrRose, thick, fill=nrPaleOrange, align=center, inner sep=1.5pt, font=\footnotesize}}
\node[nrbox, text width=2.55cm] (s1) {1. 状态采集\\$x_{i,k},q_{L,k}$\\$u_{k-1}$\\FDI 残差};
\node[nrmodule, text width=2.65cm, right=of s1] (s2) {2. 网络质量更新\\$q_{ij,k},\tau_{ij,k},\ell_{ij,k}$\\拓扑 $\mathcal G_k$\\EWMA/阈值};
\node[nrmodule, text width=2.65cm, right=of s2] (s3) {3. 邻居状态补偿\\可通信：延迟外推\\丢包：保持/短预测\\$\hat x^{net}_{j,k|k}$};
\node[nrmodule, text width=2.65cm, right=of s3] (s4) {4. Koopman 预测\\lifting $z_k$\\双线性 rollout\\残差界 $\bar w(q)$};
\node[nrmodule, text width=2.75cm, right=of s4] (s5) {5. 构造 MPC\\代价：$e_L,e_y,e_\psi,e_c$\\一致性权重 $\omega(q)$\\输入平滑 $\Delta u$};

\node[nrsafe, text width=2.55cm, below=of s1] (s9) {9. 执行与记录\\施加 $u^{act}$\\记录轨迹/连接/受力\\更新日志};
\node[nrdecision, text width=2.25cm, below=of s2] (s8) {8. 证书/约束\\通过?};
\node[nrsafe, text width=2.65cm, below=of s3] (s7) {7. 安全后处理\\FDI/FTC 重分配\\角色调度\\输入投影};
\node[nrmodule, text width=2.65cm, below=of s4] (s6) {6. 求解优化\\约束：$u,\delta,a_x$\\$\|e_c\|,\|F_c\|$\\输出 $u^\star_k$};
\node[nrsafe, text width=2.75cm, below=of s5] (fallback) {保护分支\\触发：低 $q$、高残差、证书负裕度\\收缩约束\\degraded fallback};

\draw[nrarr] (s1) -- (s2);
\draw[nrarr] (s2) -- (s3);
\draw[nrarr] (s3) -- (s4);
\draw[nrarr] (s4) -- (s5);
\draw[nrarr] (s5) -- (s6);
\draw[nrarr] (s6) -- (s7);
\draw[nrarr] (s7) -- (s8);
\draw[nrarr] (s8) -- node[above, font=\scriptsize, color=nrInk]{是} (s9);
\draw[nrdash] (s8) -- node[above, font=\scriptsize, color=nrRose]{否} (fallback);
\draw[nrarr] (fallback) -- (s6);
\draw[nrdash] (s9.west) -- ++(-0.45,0) |- node[pos=0.20, left, font=\scriptsize, color=nrRose]{下一采样周期 $k\!+\!1$} (s1.west);
\draw[nrdash] (s2.south) -- node[right, font=\scriptsize, color=nrRose]{低链路质量} (fallback.north west);
\draw[nrdash] (s1.south) -- node[left, font=\scriptsize, color=nrRose]{故障残差异常} (s7.west);
\end{tikzpicture}}
\caption{\Method{} 详细在线控制流程图。含义：该图给出每个采样周期从状态采集到控制执行的完整计算链。步骤 1--3 将车辆/载荷状态和网络质量转化为可用于优化的邻居预测；步骤 4--6 使用双线性 Koopman rollout 构造并求解一致性 MPC；步骤 7--8 对求解结果进行 FDI/FTC 重分配、角色调度、约束检查和 Lyapunov/ISS 证书检查；若通信质量、故障残差或证书裕度不满足要求，则进入保护分支并重新收缩约束/降级求解；步骤 9 才施加控制并记录实验审计指标。}
\label{fig:nrkdcc_flowchart}
\end{figure*}

\subsection{离线数据生成}

离线数据由式 \eqref{eq:world_kinematics}--\eqref{eq:connection_force} 的仿真模型生成，覆盖 nominal clean、high communication noise、single fault 和 mixed fault/noise 四类工况。训练数据采用 $\Delta t=0.02$ s，单车状态维数 $n_x=6$，输入维数 $n_u=2$；参考速度从 $2.2$--$4.0$ m/s 分层采样，初始路径进度随机平移 $0$--$80$ m，横向误差初值在 $[-2,2]$ m 内扰动。每个 batch 生成 36 条轨迹和约 1400 个 snapshot，并按 80/20 划分训练集与验证集。每个 run 同步记录车辆状态、载荷中心、连接误差、控制输入、通信质量、时延/丢包、故障效率和证书项。为避免训练数据只覆盖平稳轨迹，输入激励包含小幅随机转角、纵向加速度扰动、曲率变化和通信质量扰动。

\subsection{双线性 Koopman 提升}

令原始状态为 $x_k\in\R^{n_x}$，输入为 $u_k\in\R^{n_u}$。Koopman lifting 定义为
\begin{equation}
z_k=\Phi_\theta(x_k)=\begin{bmatrix}x_k\\ \phi_\theta(x_k)\end{bmatrix}\in\R^{n_z},
\label{eq:koopman_lift}
\end{equation}
其中 $\phi_\theta$ 为神经网络特征。本文采用双线性 lifted dynamics 表示车辆输入与载荷误差之间的乘性耦合：
\begin{equation}
z_{k+1}=A z_k+B u_k+\sum_{\ell=1}^{n_u} u_{k,\ell} N_\ell z_k+w_k,
\label{eq:bilinear_koopman}
\end{equation}
其中 $A,B,N_\ell$ 由离线数据拟合，$w_k$ 为残差。该结构比纯线性 Koopman 模型更适合描述转角、纵向加速度与车辆侧向/载荷误差之间的乘性耦合。

给定数据集 $\mathcal D=\{z_k,u_k,z_{k+1}\}$，定义回归矩阵
\begin{equation}
\Xi_k=\begin{bmatrix}z_k^\top&u_k^\top&(u_{k,1}z_k)^\top&\cdots&(u_{k,n_u}z_k)^\top\end{bmatrix}^\top .
\label{eq:bilinear_regressor}
\end{equation}
参数矩阵 $\Theta=[A\;B\;N_1\;\cdots N_{n_u}]$ 由岭回归求解：
\begin{equation}
\Theta^\star=Z_+\Xi^\top(\Xi\Xi^\top+\lambda I)^{-1}.
\label{eq:ridge_fit}
\end{equation}
其中 $\lambda$ 为岭回归正则系数，用于抑制小样本、多重共线和噪声导致的参数放大。在线阶段采用滑动窗口残差更新，只允许在稳定投影和约束边界内微调 $\Theta$，避免自适应过程破坏闭环可行性。

\subsection{稳定投影与残差边界}

为限制长 horizon 预测漂移，对 Koopman 主矩阵执行谱投影：
\begin{equation}
A\leftarrow \Pi_\rho(A),\qquad
\rho(\Pi_\rho(A))\le \bar\rho<1+\epsilon_\rho ,
\label{eq:spectral_projection}
\end{equation}
其中 $\rho(\cdot)$ 为谱半径。投影并不保证真实非线性系统全局渐近稳定，而是在 MPC 工作区域内限制 lifted predictor 的误差放大。令一步残差满足
\begin{equation}
\norm{w_k}\le \bar w_0+\bar w_x\norm{x_k-x_k^\mathrm{ref}}+\bar w_q(1-q_k),
\label{eq:residual_bound}
\end{equation}
则通信质量下降会通过 $\bar w_q(1-q_k)$ 进入后续 Lyapunov 证明中的扰动项。

\subsection{通信质量感知一致性 MPC}

MPC 在每个采样时刻求解 horizon $H$ 内的输入序列。代价函数为
\begin{equation}
\begin{aligned}
J=&\sum_{h=0}^{H-1}
\Big[
\norm{e_{L,k+h}}_{Q_L}^2+
\sum_i\norm{e_{y,i,k+h},e_{\psi,i,k+h}}_{Q_v}^2\\
&+\sum_{(j,i)\in\mathcal E_k}\omega_{ij,k}\norm{x_{i,k+h}-\hat x_{j,k+h|k}^{\mathrm{net}}}_{Q_c}^2\\
&+\norm{e_{c,k+h}}_{Q_e}^2+
\norm{u_{k+h}}_{R_u}^2+
\norm{\Delta u_{k+h}}_{R_\Delta}^2
\Big]\\
&+\norm{e_{L,k+H}}_{P_L}^2 ,
\end{aligned}
\label{eq:mpc_cost}
\end{equation}
其中 $e_L=[X_L-X_L^{\mathrm{ref}},Y_L-Y_L^{\mathrm{ref}},\psi_L-\psi_L^{\mathrm{ref}}]^\top$。通信权重定义为
\begin{equation}
\omega_{ij,k}=\omega_{\min}+(\omega_{\max}-\omega_{\min})\bar q_{ij,k},
\label{eq:comm_weight}
\end{equation}
$\bar q_{ij,k}$ 为平滑后的链路质量。控制和安全约束包括
\begin{equation}
\begin{gathered}
u_{\min}\le u_i\le u_{\max},\quad
\delta_{\min}\le\delta_i\le\delta_{\max},\quad
a_{\min}\le a_{x,i}\le a_{\max},\\
\norm{e_{c,i}}\le e_{c,\max}(q_k),\qquad
\norm{F_{c,i}}\le F_{c,\max},\\
\norm{\Delta \delta_i}\le \Delta\delta_{\max},\quad
\norm{\Delta a_{x,i}}\le \Delta a_{\max}.
\end{gathered}
\label{eq:constraints}
\end{equation}
当 $q_k$ 下降时，$e_{c,\max}(q_k)$ 与输入边界按比例收缩；当链路质量低于阈值时，degraded fallback 降低不可靠邻居权重并优先满足连接误差与输入约束。

\subsection{FDI/FTC 与相位角色调度}

执行器效率记为 $\eta_{i,k}\in[\eta_{\min},1]$，实际输入近似为
\begin{equation}
u_{i,k}^{\mathrm{act}}=\eta_{i,k}u_{i,k}^{\mathrm{cmd}}+\nu_{i,k}.
\label{eq:actuator_efficiency}
\end{equation}
FDI 监测器基于预测残差和控制响应估计 $\hat\eta_{i,k}$。设 $r_{i,k}$ 为车辆 $i$ 的一步预测残差，在线监测量采用 EWMA 与 CUSUM 组合：
\begin{equation}
\bar r_{i,k}=(1-\alpha_r)\bar r_{i,k-1}+\alpha_r\norm{r_{i,k}},\qquad
c_{i,k}=\max\{0,c_{i,k-1}+\bar r_{i,k}-r_0\}.
\label{eq:fdi_ewma_cusum}
\end{equation}
实现中使用预热步数、连续触发步数和释放步数过滤瞬时噪声；当前配置中残差阈值约为 $0.18$，CUSUM 阈值约为 $0.30$，连续检测保持为 5 步，释放保持为 10 步。执行器效率估计分别跟踪纵向加速度通道和转角通道，若 $\hat\eta^a_i<0.72$ 或 $\hat\eta^\delta_i<0.74$，则判定对应通道退化；若估计低于严重退化阈值，则触发更保守的 fallback。该设计把故障判断从单步残差变成“残差累计 + 通道效率”的联合证据，降低通信噪声造成的误报。

当 $\hat\eta_{i,k}$ 降低或残差超过阈值时，FTC 模块通过输入约束重分配和模式切换调整车辆贡献：
\begin{equation}
u_k^\mathrm{safe}=\arg\min_{u\in\mathcal U(\hat\eta,q)}
\norm{u-u_k^\mathrm{mpc}}_{R_s}^2+
\alpha_e\norm{e_{c,k+1}(u)}^2 .
\label{eq:ftc_projection}
\end{equation}
相位角色调度器根据路径曲率、通信质量和故障状态调整不同车辆在横向纠偏、纵向推进和连接保护中的权重。当前结果显示 phase-role 对部分指标有小幅收益，但不是所有指标单调改善，因此主文把它写为调度增强，而不是单独核心理论贡献。

\section{稳定性证明}

本节给出与 \Method{} 模块一致的实践稳定性证明。证明目标不是全局渐近稳定，而是在有界模型误差、通信退化可恢复和 MPC 可行的条件下，闭环跟踪误差、连接误差和 lifted 预测误差一致最终有界。

\begin{assumption}[局部 Lipschitz 与有界工作域]
闭环状态位于紧集 $\mathcal X$，输入位于紧集 $\mathcal U$。车辆--载荷真实动力学 $f$、Koopman lifting $\Phi_\theta$ 和投影 $\Pi_x$ 在该集合内局部 Lipschitz。
\end{assumption}

\begin{assumption}[预测误差与通信退化]
Koopman 一步残差满足式 \eqref{eq:residual_bound}。通信质量 $q_k$、平均时延 $\tau_k$ 和丢包率有界；严重通信退化不要求永久消失，但其连续持续时间有上界，或由 degraded fallback 保证可行输入集合非空。丢包/时延导致的邻居预测误差满足 $\norm{\tilde x_{j,k}^{\mathrm{net}}}\le \bar d_q(1-q_k)+\bar d_\tau \tau_k$。
\end{assumption}

\begin{assumption}[局部可行与证书触发保护]
在 nominal 工作域内，存在局部控制律 $\kappa_f$ 和终端集 $\mathcal X_f$，使得收缩后的输入、连接误差和载荷受力约束在有限 horizon 内递归可行。若在线求解、FDI/FTC 切换或通信退化使候选序列不可行，则保护分支只要求保持一组保守安全输入非空，并允许收缩参考推进速度。这一假设弱于全局递归可行，只覆盖本文仿真工作域。
\end{assumption}

定义组合误差
\begin{equation}
\xi_k=\mathrm{col}\left(e_{L,k},e_{c,k},e_{y,1,k},e_{\psi,1,k},\ldots,e_{y,N,k},e_{\psi,N,k},\tilde z_k\right),
\label{eq:xi_def}
\end{equation}
其中 $\tilde z_k=z_k-z_k^\star$ 为 lifted 预测误差。选取 Lyapunov 函数
\begin{equation}
V_k=\xi_k^\top P\xi_k+\sum_i \beta_i(1-\hat\eta_{i,k})^2+\gamma\sum_{(j,i)\in\mathcal E_k}(1-\bar q_{ij,k})^2 ,
\label{eq:lyapunov}
\end{equation}
$P\succ0$，$\beta_i,\gamma>0$。第一项度量跟踪、连接和预测误差；第二项度量执行器效率退化；第三项度量网络退化造成的调度风险。

在线证书采用分模式矩阵 $P_{\sigma_k}$ 和衰减率 $\lambda_{\sigma_k}$ 记录一步裕度：
\begin{equation}
m_k=(1-\lambda_{\sigma_k}\Delta t)V_{\sigma_k,k}
+\gamma_d\norm{d_k}^2+\gamma_w\norm{w_k}^2
-V_{\sigma_{k+1},k+1}.
\label{eq:mode_certificate_margin}
\end{equation}
其中 $\sigma_k\in\{\mathrm{nominal},\mathrm{FTC},\mathrm{degraded}\}$。若 $m_k\ge -\epsilon_m$，则该步证书通过；若连续负裕度出现，控制器降低推进、收缩连接约束或切入 degraded fallback。后续实验中的 ok ratio 即统计式 \eqref{eq:mode_certificate_margin} 的通过比例。

\begin{theorem}[输入到状态实践有界]
若假设 1--3 成立，并且 MPC 代价权重满足
\begin{equation}
Q_L,Q_e,Q_v,R_u,R_\Delta\succ0,
\label{eq:positive_weights}
\end{equation}
稳定投影满足 $\rho(\Pi_\rho(A))\le \bar\rho<1+\epsilon_\rho$，且分模式 MPC 与保护分支在证书通过步满足存在 $\alpha\in(0,1)$、$c_d>0$ 使得
\begin{equation}
V_{k+1}-V_k\le -\alpha\norm{\xi_k}^2+c_d\norm{d_k}^2+c_q(1-q_k)^2+c_\tau\tau_k^2,
\label{eq:iss_difference}
\end{equation}
并且负裕度步数满足有限密度约束，即存在 $\bar \nu<1$ 使得任意长度窗口内证书失败比例不超过 $\bar \nu$，则闭环误差 $\xi_k$ 对扰动 $d_k$、通信质量退化 $1-q_k$ 和时延 $\tau_k$ 输入到状态实践有界。特别地，若这些外部项一致有界，则存在常数 $C>0$、$\lambda\in(0,1)$ 和最终界 $\Omega$，使得
\begin{equation}
\norm{\xi_k}^2\le C\lambda^k\norm{\xi_0}^2+\Omega ,
\label{eq:uub_bound}
\end{equation}
其中 $\Omega=O(\bar d^2+\overline{(1-q)}^2+\bar\tau^2+\bar w_0^2)$。
\end{theorem}

\begin{proof}
由 MPC 最优性，时刻 $k$ 的最优序列去掉首项并接上终端控制律 $\kappa_f$ 可构造时刻 $k+1$ 的候选可行序列。递归可行假设保证该候选序列满足收缩后的输入、连接和受力约束。因此 nominal 情况下，标准 MPC 下降性质给出
\begin{equation}
J_{k+1}^\star-J_k^\star\le -\ell(\xi_k,u_k)+c_w\norm{w_k}^2 ,
\label{eq:mpc_descent}
\end{equation}
其中 $\ell$ 为阶段代价下界，满足 $\ell(\xi_k,u_k)\ge \underline q\norm{\xi_k}^2$。双线性 Koopman 残差、邻居预测误差和执行器效率估计误差会使候选轨迹与真实轨迹偏离；由局部 Lipschitz 性可将这些偏离上界写为
\begin{equation}
\norm{\Delta \xi_{k+1}}\le L_w\norm{w_k}+L_q(1-q_k)+L_\tau\tau_k+L_d\norm{d_k}.
\label{eq:error_perturbation}
\end{equation}
将式 \eqref{eq:residual_bound} 与式 \eqref{eq:error_perturbation} 代入式 \eqref{eq:mpc_descent}，并用 Young 不等式吸收交叉项，可得
\begin{equation}
J_{k+1}^\star-J_k^\star\le -\alpha_0\norm{\xi_k}^2+c_d\norm{d_k}^2+c_q(1-q_k)^2+c_\tau\tau_k^2+c_0 .
\label{eq:descent_with_disturbance}
\end{equation}
FDI/FTC 项在无故障时不增加 $V_k$；在故障或通信退化触发时，控制重分配求解式 \eqref{eq:ftc_projection}，保证输入落入保守可行集，并优先减小连接误差项。有限密度负裕度约束排除了证书长期连续失败，因而切换和 fallback 带来的 Lyapunov 增量可由有限常数界定并并入 $c_0$。令 $V_k$ 与 $J_k^\star$ 在工作域内等价，即 $\underline c\norm{\xi_k}^2\le V_k\le \bar c\norm{\xi_k}^2+\bar c_0$，得到式 \eqref{eq:iss_difference}。对该差分不等式迭代求和即可得到式 \eqref{eq:uub_bound}。
\end{proof}

\begin{remark}
该证明与当前实验中的证书记录一致：E3 只能作为证书裕度和 ok ratio 的闭环证据，不能替代理论证明，也不能推出所有场景全局渐近稳定。旧版全场景聚合统计中 nominal 和 single-fault 出现负最小裕度；本轮补充的分模式压力场景证据显示 high-communication-noise 与 mixed fault/noise 中 nominal/FTC 模式均为正裕度。因此主文只写“分模式实践有界与安全监测”，不写无条件渐近稳定。
\end{remark}

\section{实验结果与分析}

\subsection{实验组与基线边界}

本文当前使用五组实验支撑方法模块。E1 为 Koopman 训练损失、控制窗口预测与稳定投影审计；E0 为主对比，比较 \AKE{}、\Method{} w/o role schedule、\Method{} 和 ZOH-consensus surrogate；E2 为模块消融，移除通信感知、时延补偿、稳定投影、FDI/FTC、phase-role、PPC guard 等模块；E3 记录 Lyapunov/ISS 证书统计；E7 在相同扰动下检查载荷连接安全和受力代价。E0/E2/E3/E7 当前已有 paired $n=20$ 结果；E1 当前支撑“训练收敛、短窗口预测可比和稳定投影有效”，不支撑“远期开环 rollout 全面优于线性模型”。

需要明确两类基线。第一，\AKE{} 是对上传基线论文思想的任务对齐机制实现，用于当前四车载荷运输场景下的同场对比；它不是对上传 AKE-A 原文全设置的严格复现。第二，ZOH-consensus surrogate 是低阶非 Koopman 辅助地板基线，只用于证明“无学习预测/无网络感知保护”的退化风险，不能替代强物理 DMPC 或鲁棒控制基线。

\subsection{Koopman、时延与 FDI/FTC 机制证据}

\begin{figure*}[!t]
\centering
\includegraphics[width=\textwidth]{figures/fig_nrkdcc_koopman_training_loss.png}
\caption{Koopman lifting 网络训练/验证损失。含义：左图给出 total、prediction 和 lifted loss 在训练集与验证集上的收敛过程；右图给出最后 10 个 epoch 的训练/验证均值差距。该图用于证明网络训练过程是可复查、可收敛且无明显发散的，而不是直接证明长时开环 rollout 优势。}
\label{fig:e1_koopman_training_loss}
\end{figure*}

图 \ref{fig:e1_koopman_training_loss} 的客观结果是，total、prediction 和 lifted 三类损失在训练集与验证集上均快速下降，并在约 60 个 epoch 后进入稳定平台；最后 10 个 epoch 的验证损失与训练损失处于同一量级，未出现明显持续发散或训练--验证分离。

造成上述现象的原因是，离线数据覆盖了横向误差、曲率和输入扰动的主要工作域，lifting 网络首先学习状态重构和一步预测的低维一致结构，随后双线性回归在 lifted 空间中拟合控制输入对状态演化的影响。因此该图能支撑“训练过程收敛且可复查”，但不能单独证明所有 horizon 上均优于线性模型；本文据此把 Koopman 模块限定为“提供可训练、可投影、可残差审计的预测结构”。

\begin{figure*}[!t]
\centering
\includegraphics[width=\textwidth]{figures/fig_nrkdcc_koopman_prediction_evidence.png}
\caption{E1 Koopman 预测与稳定投影审计。含义：左图给出不同 lifted predictor 的一步状态 RMSE；中图只展示与当前 MPC 预测时域一致的 1--12 步 rollout 误差及 95\% 置信区间；右图给出稳定投影前后的谱半径。该图支持“短窗口预测不劣于线性近似、稳定投影把谱半径压到 1 以下、预测器可被在线残差审计”的机制结论。}
\label{fig:e1_koopman_evidence}
\end{figure*}

图 \ref{fig:e1_koopman_evidence} 的客观结果是，双线性模型的一步平均 RMSE 为 $0.0339$，略低于线性模型的 $0.0350$；在 1--8 步短窗口内，双线性与稳定双线性 rollout 平均误差分别约为 $0.2716$ 和 $0.2763$，低于线性模型的 $0.2797$；到当前 MPC 主配置的 $H=12$ 时，三类模型的 95\% 置信区间高度重叠。右图显示稳定投影后谱半径由 $1.006$ 降至 $0.999$。

造成上述结果的原因是，双线性 Koopman 结构能在短预测窗口内表达输入--状态耦合，因此一步和短窗口误差略有收益；但开放环 rollout 会逐步积累模型残差、状态归一化误差和曲率外推误差，所以到 $H=12$ 后优势不再显著。谱投影的作用不是直接降低所有 horizon 的平均误差，而是限制不稳定特征值导致的误差放大，因此本文只能写“控制窗口内可比且短窗口略优”，不能写“长时开环全面优于线性模型”。

\begin{figure*}[!t]
\centering
\includegraphics[width=\textwidth]{figures/fig_nrkdcc_delay_mechanism_evidence.png}
\caption{通信时延补偿机制与消融边界。含义：左图显示 mixed fault/noise 中在线平均时延、丢包率、链路质量和约束收缩随时间变化；右图显示 high-communication-noise 场景下 full \Method{} 与 w/o delay 的 paired $n=20$ 指标。该图说明时延补偿模块已接入通信质量估计和约束收缩链条，但当前时延强度下 w/o delay 与 full 方法差异很小，因此时延补偿只可写为网络韧性机制，不能写为已充分证明的主要性能增益。}
\label{fig:delay_mechanism_evidence}
\end{figure*}

图 \ref{fig:delay_mechanism_evidence} 的客观结果是，mixed fault/noise 工况下平均时延、丢包率和链路质量随时间变化，控制器同步给出非零约束收缩；在 high-communication-noise 消融中，移除 delay compensation 后横向 RMSE、纵向 RMSE、连接利用率和证书 ok ratio 与 full \Method{} 接近，差异并不显著。

造成上述结果的原因是，当前通信扰动强度下，主导闭环安全的是链路质量权重和约束收缩，而不是单独的时延外推项；当延迟幅值和连续丢包长度不足以显著破坏邻居预测时，no-delay-compensation 与 full 方法会表现接近。因此本文把时延补偿写成网络韧性机制和可解释保护链条，而不把它写成已经充分证明的主要性能增益。

\begin{figure*}[!t]
\centering
\includegraphics[width=\textwidth]{figures/fig_nrkdcc_fdi_ftc_timeline.png}
\caption{FDI/FTC 逐步诊断与重构时间线。含义：该图在 mixed fault/noise 单次可复查 run 中标出故障注入、FDI 诊断和 FTC 模式切换时刻，同时显示故障缺口、诊断置信度、重分配控制量和证书裕度。故障约在 $7.26$ s 激活，FDI 约 10 个采样步后给出非名义诊断，FTC 约 13 个采样步后切换到 reconfigured mode。}
\label{fig:fdi_ftc_timeline}
\end{figure*}

图 \ref{fig:fdi_ftc_timeline} 的客观结果是，故障约在 $7.26$ s 激活后，残差缺口上升，FDI 置信度从零进入退化判定区；随后全局模式由 nominal Koopman MPC 切换为 reconfigured FTC MPC，并出现非零纵向加速度和转角重分配量。证书裕度在切换过程中保持可监测状态。

造成上述结果的原因是，执行器效率下降首先表现为期望输入与实际响应之间的残差缺口，FDI 模块据此累积置信度并触发故障模式；FTC 模块随后通过保守可行集内的输入重分配补偿退化车辆，避免连接误差继续放大。该图能证明 FDI/FTC 已接入闭环并按预期触发，但由于它是单次时间线，不能单独宣称所有故障类型均可快速精确识别。

\begin{figure*}[!t]
\centering
\includegraphics[width=\textwidth]{figures/fig_nrkdcc_modewise_certificate_closure.png}
\caption{E3 分模式 Lyapunov/ISS 证书闭合结果。含义：该图把 high-communication-noise 与 mixed fault/noise 中 nominal/FTC 模式分开统计，展示证书 ok ratio、最小裕度和 95\% 实践界。分模式统计均为正裕度，说明压力场景下证书与保护逻辑一致；这修正了旧版全场景聚合统计中 nominal/single-fault 负裕度导致的过强稳定性表述。}
\label{fig:modewise_certificate_closure}
\end{figure*}

图 \ref{fig:modewise_certificate_closure} 的客观结果是，high-communication-noise 的 nominal 模式、mixed fault/noise 的 nominal 模式以及 reconfigured FTC 模式均达到 ok ratio $1.0$，且最小裕度为正；与此相对，旧版全场景表 \ref{tab:e3_certificate} 中 nominal clean 与 single fault 的最小裕度存在负值。

造成上述差异的原因是，全场景最小值会把短时数值扰动、切换瞬态和非压力工况中的局部负裕度统一压缩成最保守指标，而分模式统计更接近实际控制逻辑中的 nominal/FTC 分支。因此本文稳定性表述应从“所有场景全局稳定”收缩为“在通信/故障压力场景下，分模式证书与 ISS/UUB 保护逻辑一致”。

\subsection{主对比结果}

\begin{table*}[!t]
\centering
\caption{E0 主对比：mixed fault/noise 与 high-communication-noise 场景下的 paired $n=20$ 统计。含义：该表给出主方法、任务对齐基线和低阶辅助基线在成功率、跟踪误差、连接安全和力峰值上的总体差异，用于界定本文最核心的性能 claim 与受力代价边界。}
\label{tab:e0_main}
\footnotesize
\begin{tabular}{llcccccc}
\toprule
场景 & 方法 & 成功率 & 全路径成功 & 横向 RMSE & 纵向 RMSE & 连接最大利用率 & 力峰值 \\
\midrule
\multirow{4}{*}{mixed fault/noise}
& \AKE{} & 1.00 & 1.00 & 0.0453 & 0.4724 & 0.2999 & 1324.7\\
& \Method{} w/o role schedule & 1.00 & 1.00 & 0.0412 & 0.4681 & 0.0497 & 2138.5\\
& \TF{} & 1.00 & 1.00 & 0.0408 & 0.4680 & 0.0513 & 2219.4\\
& ZOH surrogate & 0.00 & 0.00 & 0.0879 & 14.0477 & 0.3248 & 1023.1\\
\midrule
\multirow{4}{*}{high comm. noise}
& \AKE{} & 1.00 & 1.00 & 0.0519 & 0.5114 & 0.3757 & 1438.5\\
& \Method{} w/o role schedule & 1.00 & 1.00 & 0.0420 & 0.4708 & 0.0373 & 2185.5\\
& \TF{} & 1.00 & 1.00 & 0.0424 & 0.4708 & 0.0452 & 2219.4\\
& ZOH surrogate & 0.00 & 0.00 & 0.1096 & 14.1466 & 0.3556 & 1775.3\\
\bottomrule
\end{tabular}
\begin{tablenotes}
\footnotesize
\item 注：力峰值越大表示代价更高，不作为本文的单调改善 claim。当前结果支持连接误差和跟踪误差改善，但同时提示受力代价需要在后续权重整定和约束设计中进一步控制。
\end{tablenotes}
\vspace{1.5mm}
\begin{minipage}{0.93\textwidth}
\footnotesize
\textbf{表内解读：} 该表不把“成功率”作为主要优势，因为 \AKE{} 在两个压力场景下同样全路径成功。本文能稳妥书写的主结论是：\Method{} 在保持任务完成的同时显著降低连接最大利用率，并同步降低横向 RMSE；但载荷力峰值升高，说明网络安全保护和受力平滑之间仍存在权衡，后文 E7 将单独审计该负面边界。
\end{minipage}
\end{table*}

表 \ref{tab:e0_main} 的客观结果是，\AKE{} 和 \Method{} 在 mixed fault/noise 与 high-communication-noise 两个压力场景下均达到全路径成功，因此成功率不是主要差异。差异主要集中在跟踪误差和连接安全：mixed fault/noise 场景中，\Method{} 相对 \AKE{} 将横向 RMSE 从 0.0453 降至 0.0408，连接最大利用率从 0.2999 降至 0.0513；high-communication-noise 场景中，横向 RMSE 从 0.0519 降至 0.0424，连接最大利用率从 0.3757 降至 0.0452。同时，\Method{} 的力峰值高于 \AKE{}。

造成上述结果的原因是，\Method{} 将通信质量权重、连接几何约束和 Koopman-MPC 预测统一进入优化问题，使不可靠邻居和高风险连接在代价与约束中被主动抑制，因此连接风险暴露和横向误差更容易下降。力峰值升高则说明控制器在当前权重下更倾向于牺牲瞬时载荷受力来维持连接几何安全，因此本文不能写“受力更小”或“载荷受力更平滑”，只能写“连接最大利用率显著下降，但受力代价需要进一步调权”。

\begin{figure*}[!t]
\centering
\includegraphics[width=0.92\textwidth]{figures/fig_nrkdcc_team_center_trajectory.png}
\caption{团队中心轨迹对比。含义：该图直接展示 \Method{}、无角色调度版本和 \AKE{} 在通信强退化与混合故障/噪声下对载荷中心参考轨迹的贴合程度，用于补充表 \ref{tab:e0_main} 中 RMSE 统计的空间直观证据。}
\label{fig:nrkdcc_team_center_trajectory}
\end{figure*}

\begin{figure*}[!t]
\centering
\includegraphics[width=0.92\textwidth]{figures/fig_nrkdcc_error_diagnostics.png}
\caption{团队中心与各车轨迹误差诊断。含义：上两排分别给出载荷中心横向误差 $e_y$ 和纵向误差 $e_s$，用于判断 \Method{} 的优势是否贯穿时间过程；下排给出 mixed fault/noise 场景下 \AKE{} 与 \Method{} 的四辆车横向误差，用于检查团队中心改善是否由单车过度补偿造成。}
\label{fig:nrkdcc_error_diagnostics}
\end{figure*}

图 \ref{fig:nrkdcc_team_center_trajectory} 的客观结果是，在 high-communication-noise 和 mixed fault/noise 两类压力场景中，\Method{} 的团队中心轨迹与参考轨迹基本保持同相位贴合。图 \ref{fig:nrkdcc_error_diagnostics} 进一步显示，0--20 s 内 \Method{} 的团队中心横向偏差峰值整体较低，纵向误差仍随正弦参考周期呈现振荡；下排单车横向误差未出现某一车辆在故障窗口后持续失控式放大。

造成上述结果的原因是，团队中心误差由载荷几何约束、各车一致性项和参考轨迹跟踪项共同决定，通信质量感知权重降低了退化链路对团队中心的错误牵引，phase-role 调度则避免单车长期承担过高修正量。纵向误差仍有周期振荡，是因为当前控制更优先处理横向连接安全和载荷中心几何约束，纵向速度/推进方向还不是当前权重配置下的最强优化目标。

\FloatBarrier

\subsection{模块消融}

\begin{table*}[!t]
\centering
\caption{E2 mixed fault/noise 全量模块消融 paired $n=20$ 统计。含义：该表逐项移除关键模块，量化各模块对任务成功、跟踪误差、连接利用率和证书 ok ratio 的贡献，用于避免只用主对比推断模块有效性。}
\label{tab:e2_mixed}
\footnotesize
\begin{tabular}{lccccc}
\toprule
方法 & 成功率 & 横向 RMSE & 纵向 RMSE & 连接最大利用率 & 证书 ok 均值 \\
\midrule
full \Method{} & 1.00 & 0.0408 & 0.4679 & 0.0550 & 0.99995\\
w/o comm-aware & 1.00 & 0.0414 & 0.5068 & 0.4294 & 0.78555\\
w/o delay comp. & 1.00 & 0.0408 & 0.4678 & 0.0556 & 0.99995\\
w/o stable proj. & 1.00 & 0.0388 & 0.4677 & 0.0560 & 0.99995\\
w/o PPC guard & 0.00 & 0.0660 & 12.9301 & 0.1259 & 0.28722\\
ZOH surrogate & 0.00 & 0.0970 & 14.0191 & 0.3324 & 0.17739\\
\bottomrule
\end{tabular}
\end{table*}

\begin{table}[!t]
\centering
\caption{E2 high-communication-noise 关键模块消融 paired $n=20$ 统计。含义：该表聚焦网络退化场景，检验通信感知与时延补偿两个网络模块对连接安全和轨迹误差的独立作用。}
\label{tab:e2_comm}
\footnotesize
\begin{tabular}{lcccc}
\toprule
方法 & 成功率 & 横向 RMSE & 纵向 RMSE & 连接利用率\\
\midrule
full \Method{} & 1.00 & 0.0425 & 0.4709 & 0.0446\\
w/o comm-aware & 1.00 & 0.0459 & 0.5399 & 0.4384\\
w/o delay comp. & 1.00 & 0.0426 & 0.4709 & 0.0460\\
ZOH surrogate & 0.00 & 0.1037 & 14.1242 & 0.3649\\
\bottomrule
\end{tabular}
\end{table}

表 \ref{tab:e2_mixed} 和表 \ref{tab:e2_comm} 的客观结果是，通信感知模块移除后退化最明显：mixed fault/noise 场景中连接最大利用率由 0.0550 升至 0.4294，证书 ok 均值由 0.99995 降至 0.78555；high-communication-noise 场景中连接利用率由 0.0446 升至 0.4384，纵向 RMSE 由 0.4709 升至 0.5399。移除 PPC guard 后，mixed fault/noise 场景全路径成功率降为 0，纵向 RMSE 上升到 12.9301；ZOH surrogate 在两个压力场景中均不能完成全路径。no-delay-compensation 与 full \Method{} 的差距较小。

造成上述结果的原因是，链路质量并非附加诊断量，而是直接改变一致性权重、邻居预测可信度和约束收缩幅度；去掉 comm-aware 后，不可靠邻居仍被等权纳入协同优化，连接误差更容易积累。PPC guard 是末端可行域保护，移除后网络退化和故障扰动会被积累成不可恢复的纵向偏移。相比之下，当前时延强度不足以单独主导退化，因此 no-delay-compensation 与 full \Method{} 表现接近。

\begin{figure}[!t]
\centering
\includegraphics[width=\linewidth]{figures/fig_nrkdcc_mixed_ablation_summary.png}
\caption{E2 mixed fault/noise 全量模块消融图。含义：该图把表 \ref{tab:e2_mixed} 中的横向误差、纵向误差、连接利用率和证书通过率画成同尺度面板，用于快速定位通信感知、PPC guard 和 ZOH surrogate 对闭环退化的影响。}
\label{fig:r2_mixed_heatmap}
\end{figure}

\begin{figure}[!t]
\centering
\includegraphics[width=\linewidth]{figures/fig_nrkdcc_comm_ablation_summary.png}
\caption{E2 high-communication-noise 关键模块消融图。含义：该图突出通信感知模块对连接利用率和纵向误差的影响，同时显示 no-delay-compensation 与 full \Method{} 差别较小，因此时延补偿仍需 dose-response 进一步证明。}
\label{fig:r2_mixed_effects}
\end{figure}

图 \ref{fig:r2_mixed_heatmap} 和图 \ref{fig:r2_mixed_effects} 的客观结果是，mixed fault/noise 图中 w/o comm-aware 的连接利用率柱显著升高，w/o PPC guard 和 ZOH 的纵向 RMSE 明显失控；high-communication-noise 图中，w/o comm-aware 的连接利用率接近 0.44，而 full \Method{} 保持在 0.05 以下。

造成上述图形差异的原因是，通信感知模块降低不可靠邻居在一致性项中的影响，并通过约束收缩避免载荷连接误差被放大；PPC guard 则提供最后的可行域保护，使扰动积累到约束边界前被截断。ZOH surrogate 缺少学习预测和网络感知保护，因此在持续通信退化下无法主动修正纵向累积误差。

\FloatBarrier

\subsection{稳定证书与连接安全}

\begin{table}[!t]
\centering
\caption{E3 证书统计 paired $n=20$。含义：该表报告 Lyapunov/ISS 证书的通过率和最小裕度，用于区分“闭环可监测实践有界”和“严格全局稳定”这两个不同层级的结论。}
\label{tab:e3_certificate}
\footnotesize
\begin{tabular}{lccc}
\toprule
场景 & 全路径成功 & 证书 ok & 最小裕度\\
\midrule
nominal clean & 1.00 & 0.360 & -0.0224\\
high comm. noise & 1.00 & 1.000 & 0.0049\\
single fault & 1.00 & 0.799 & -0.0111\\
mixed fault/noise & 1.00 & 1.000 & 0.0033\\
\bottomrule
\end{tabular}
\end{table}

\begin{figure}[!t]
\centering
\includegraphics[width=\linewidth]{figures/fig_nrkdcc_certificate_summary.png}
\caption{E3 Lyapunov/ISS 证书统计。含义：comm-high 与 mixed 场景为正裕度，但 nominal 与 single-fault 出现负最小裕度，故该图只能支撑实践有界与监测结论，不能支撑全局渐近稳定。}
\label{fig:r3_certificate}
\end{figure}

\begin{table}[!t]
\centering
\caption{E7 连接安全与受力代价 paired $n=20$。含义：该表把连接最大利用率与力峰值并列报告，用于说明 \Method{} 的安全收益主要体现在连接约束利用率下降，同时受力代价仍需后续优化。}
\label{tab:e7_safety}
\footnotesize
\begin{tabular}{llccc}
\toprule
场景 & 方法 & 全路径成功 & 连接利用率 & 力峰值\\
\midrule
high comm. & \AKE{} & 1.00 & 0.3308 & 1315.4\\
high comm. & \TF{} & 1.00 & 0.0447 & 2212.7\\
mixed & \AKE{} & 1.00 & 0.3021 & 1326.8\\
mixed & \TF{} & 1.00 & 0.0501 & 2171.1\\
\bottomrule
\end{tabular}
\end{table}

\begin{figure}[!t]
\centering
\includegraphics[width=\linewidth]{figures/fig_nrkdcc_connection_force_audit.png}
\caption{E7 连接安全与载荷受力审计图。含义：该图把 high-communication-noise 与 mixed fault/noise 两个场景下的连接最大利用率和力峰值并列显示，说明 \Method{} 显著降低连接利用率，但力峰值更高，后续需要继续按 $F_x,F_y,M_z$ 分量调权。}
\label{fig:r7_comm}
\end{figure}

表 \ref{tab:e3_certificate} 与图 \ref{fig:r3_certificate} 的客观结果是，high-communication-noise 与 mixed fault/noise 场景的证书 ok 和最小裕度表现较好，但 nominal clean 和 single fault 的全场景最小裕度为负。图 \ref{fig:modewise_certificate_closure} 进一步按压力场景和控制模式拆分后，high-communication-noise 与 mixed fault/noise 的 nominal/FTC 模式均为正裕度。

造成上述证书差异的原因是，全场景最小裕度对短时数值误差、切换瞬态和非压力工况中的局部违背非常敏感，而分模式证书更符合实际控制器的 nominal/FTC/fallback 逻辑。因此本文只能主张可监测的分模式 ISS/UUB 证书和安全保护逻辑，稳定性 claim 必须限定在有界工作域、有限密度负裕度和 fallback 可行条件下。

表 \ref{tab:e7_safety} 与图 \ref{fig:r7_comm} 的客观结果是，\Method{} 在 high-communication-noise 场景把连接利用率从 0.3308 降至 0.0447，在 mixed fault/noise 场景从 0.3021 降至 0.0501；同时，力峰值从 1315.4/1326.8 升至 2212.7/2171.1。

造成上述安全--受力权衡的原因是，当前 MPC 权重和约束优先降低连接几何误差与载荷脱联风险，在通信退化或故障扰动下会使用更强的瞬时控制来维持队形连接，从而推高载荷受力峰值。后续实验必须补充 $F_x,F_y,M_z$ 分量图和力平滑权重消融，才能把受力指标写成正向贡献。

\FloatBarrier

\section{讨论与剩余缺口}

当前证据链最强的结论是：在四车载荷协同运输的 mixed fault/noise 和 high-communication-noise 场景下，\Method{} 相对 \AKE{} 明显降低连接最大利用率，并降低横向/纵向跟踪误差；通信感知一致性 MPC 是目前证据最充分的模块。消融中移除通信感知后，high-communication-noise 场景连接最大利用率由 0.0446 升至 0.4384，证书 ok 均值从 0.99995 降至 0.31775，说明链路质量进入一致性和约束逻辑具有明确内在机理。本轮新增的 Koopman、时延、FDI/FTC 和分模式证书图把这些模块的“是否接入闭环、何时触发、证书是否通过”补成可复查证据链。

当前不能过度书写的结论有四点。第一，\AKE{} 仍是机制对齐 baseline，不是上传原文 AKE-A 的严格复现；若要达到严格顶刊标准，需要复现 AKE-A 或补强物理/DMPC baseline。第二，ZOH surrogate 只可作为辅助地板基线。第三，force peak 在多个场景中高于 \AKE{}，因此不能写“受力更小”或“受力更平滑”。第四，no-delay-compensation 在当前 mixed 与 comm-key 消融中与 full \Method{} 差异很小，时延补偿已完成机制证据闭合，但其独立性能优势仍需要高延迟强度、不同丢包率和拓扑恢复曲线进一步验证。

下一轮应优先补四类实验：E0 all-scenarios $n=20$ 完整主对比；delay/dropout dose-response 与 topology recovery；若要扩展 Koopman claim，则补 fresh training 多 seed 和远期开环/闭环滚动预测对比；至少一个强物理/DMPC baseline 或上传 AKE-A 严格复现。同时需要补 C4 近三年 actuator FDI/FTC 直接文献，并把当前单 run 的 FDI/FTC 时间线扩展为多 seed 检测延迟、误报率和切换延迟统计。

\section{结论}

本文形成了面向网络退化四车协同运输的 \Method{} 中文 LaTeX 初稿。方法上，本文将稳定投影双线性 Koopman 预测器、通信质量感知一致性 MPC、FDI/FTC 降级和 Lyapunov/ISS 证书统一到载荷协同运输闭环中。实验上，当前 paired $n=20$ 结果支持 \Method{} 在连接误差和轨迹跟踪方面相对 \AKE{} 的改善，并通过消融证明通信感知模块是主要性能来源。本轮新增证据进一步补上 Koopman 训练损失、控制窗口预测、稳定投影审计、时延机制边界、FDI/FTC 触发时间线和分模式正裕度证书。本文同时保留严格 claim 边界：受力代价尚未单调改善，时延补偿独立性能优势仍需 dose-response 证据，严格 AKE-A/强 DMPC baseline 仍是投稿前必须补齐的缺口；若要把 Koopman claim 扩展到远期开环优势，还需要重新训练并加入多步 rollout loss 或闭环滚动预测证据。

\section*{致谢}

\todozh{补充基金项目、课题编号和作者贡献声明。}

\FloatBarrier
\balance
\bibliographystyle{elsarticle-num}
\bibliography{core_references_round10}

\end{document}
